\problema{Merge Intervals}{56}{src/03-intervalos/56-merge-intervals.cpp}{M}

Nos dan un arreglo de intervalos \texttt{intervals[i] = [inicio, fin]}.
Queremos \emph{fusionar} todos los intervalos que se solapan entre sí (o que
apenas se tocan en un punto) y devolver el arreglo resultante de intervalos
que ya no se solapan, cubriendo exactamente el mismo rango total que el
arreglo original.

\paragraph{La idea, con un ejemplo de la vida real.} Piensa en tu calendario
de reuniones de la semana, donde cada reunión ocupa un bloque de tiempo.
Si dos reuniones se traslapan (o una empieza justo cuando termina la otra),
en la práctica ocupan tu agenda como si fueran una sola reunión larga.
Fusionar intervalos es exactamente eso: convertir un montón de bloques que
se pisan entre sí en la lista más corta posible de bloques ``limpios'' que
representan lo mismo.

\begin{center}
\ttfamily
\begin{tabular}{l}
1--3\\
\ \ 2-----6 \hspace{1cm} 8--10 \hspace{1cm} 15---18\\
\end{tabular}
\end{center}

\noindent Con \texttt{intervals = [[1,3],[2,6],[8,10],[15,18]]}: \texttt{[1,3]}
y \texttt{[2,6]} se solapan (porque $3 \geq 2$), así que se funden en
\texttt{[1,6]}. Los otros dos no tocan a nadie más. El resultado es
\texttt{[[1,6],[8,10],[15,18]]}.

\paragraph{Cómo lo resuelve el código, paso a paso.}
\begin{enumerate}
  \item Ordenamos los intervalos por su valor de \texttt{inicio}. Esto es
        el paso clave: una vez ordenados, cualquier intervalo que se pueda
        solapar con uno anterior tiene que estar cerca de él en la lista, no
        disperso por todo el arreglo.
  \item Recorremos los intervalos ya ordenados manteniendo un
        ``\texttt{resultado}'' que empieza vacío.
  \item Para cada intervalo, lo comparamos contra el \emph{último} grupo que
        llevamos armado en \texttt{resultado}: si el fin de ese último grupo
        es menor que el inicio del intervalo actual, no hay solape, así que
        el intervalo actual abre un grupo nuevo.
  \item Si sí hay solape (o se tocan justo en un punto), no agregamos un
        grupo nuevo: simplemente estiramos el fin del último grupo hasta el
        mayor de los dos fines.
  \item Al terminar el recorrido, \texttt{resultado} contiene los intervalos
        fusionados, en orden.
\end{enumerate}

\lstinputlisting{src/03-intervalos/56-merge-intervals.cpp}

\complejidad{$O(n \log n)$}{$O(n)$}

\paragraph{¿Qué significa esa complejidad?} El costo dominante es ordenar
los $n$ intervalos, que toma $O(n \log n)$; el recorrido que fusiona ya
ordenados es una sola pasada, $O(n)$. En espacio, en el peor caso (ningún
intervalo se solapa con otro) el resultado tiene los mismos $n$ intervalos
que la entrada.

\paragraph{Consejos para la entrevista (Amazon).} El error más común es
olvidar ordenar primero: sin ordenar, comparar solo contra el ``último''
grupo no sirve de nada, porque el solape puede estar en cualquier parte del
arreglo. Otro punto que preguntan seguido: ¿qué pasa si dos intervalos solo
se \emph{tocan} (el fin de uno es igual al inicio del otro), se consideran
solapados? Según el enunciado de LeetCode, sí, hay que fusionarlos (usa
\texttt{<} estricto al comparar, no \texttt{<=}, para decidir cuándo NO se
solapan). Menciona también el caso de un arreglo vacío o de un solo
intervalo, que deben devolverse tal cual.
